% Root proof; D independent complete literal audit passed; acyclic excess-eight root+B gate passed.
\subsection{An excess-eight bound without a support restriction}

\begin{proposition}\label{C:strong-pair-excess-eight-through21}
Let $C$ be a connected finite category with $k$ objects and $n$ arrows,
including identities, such that
\[
 n\leq\min\{21,2k+7\}.
\]
Suppose that no object is a retract of another. If $C$ has a proper
strongly connected component on two objects, then
$BC\not\simeq\mathbb CP^2$.
\end{proposition}
\begin{proof}
Write $S$ for such a component, $s=|S|$, and $\ell=k-2\geq1$.
Both identities of $S$ lie outside its minimal ideal by the
minimal-group inverse adjustment in Lemma~\ref{C:finite-category-Rees}.
Consequently $s\geq6$. Contracting $S$ in the undirected support graph
leaves $\ell+1$ vertices. Its at least $\ell$ support edges include
an interface with $S$, which supplies arrows at both objects of $S$.
Together with the $\ell$ outside identities this gives
\[
 n\geq s+2\ell+1,\qquad n\leq2\ell+11,\qquad s\leq10.
\]
The rational categorical Rees comparison then gives
$H^2(BS;\mathbb Q)=0$: nonzero second homology requires both total
row and total column counts at least three, and therefore at least
nine minimal-ideal arrows and the two additional identities.

\paragraph{An extreme component.}
Take opposites to make an extreme $S$ a source. Let
$D=C\setminus S$ have $q$ connected components. If
$H^2(BD;\mathbb Q)\ne0$, Lemma~\ref{C:arrow-excess-two-H2}
forces at least $2t+2$ arrows in one connected component with $t$
objects. Each other component with $t'$ objects has at least
$2t'-1$ arrows. Thus $|D|\geq2\ell-q+3$.
Every component of $D$ meets $S$ in at least two cross arrows,
so
\[
 n\geq s+2\ell+q+3\geq s+2\ell+4.
\]
For $s\geq8$ this contradicts $n\leq2\ell+11$.
Both vertex blocks consequently have zero rational second
cohomology in that case, and
Lemma~\ref{C:two-block-zero-H2-products} makes every degree-two
product zero.

For $s=6$ or $7$, the minimal group is trivial and there is one
minimal row and one minimal column at each object.
Lemma~\ref{C:one-row-source-block-collapse} replaces $S$ by a point
through an actual homotopy equivalence. It saves $s-1$ internal
arrows and at least one cross arrow: for any nonempty receiving
stalk family, its orbit set is an image of either of the two
nonempty stalks, so its size is at most their smaller size.
No other hom-set increases. The resulting category has at most
$n-s\leq15$ arrows, contrary to Theorem~\ref{C:CP2-lower17}.

\paragraph{A nonextreme component.}
There are outside objects $a,b$ with $a\to S\to b$.
They are distinct, and composition supplies $a\to b$.
Thus the contracted support has at least $\ell+1$ edges,
at least two of which meet $S$. Let $P_{\rm ext}$ be the sum
of $|C(d,i)|-1$ and $|C(i,d)|-1$ over nonempty hom-sets with
$i\in S$ and $d\notin S$. Counting as above gives
\[
 n\geq s+2\ell+3+P_{\rm ext},\qquad
 P_{\rm ext}\leq8-s,\qquad s\leq8.
\]
For $s=6$ or $7$, let $e_i$ be the unique minimal-ideal
endomorphism at each of the two objects $i$. It is a zero
of its corner. Consider the four conditions asserting that
$e_i$ fixes every incoming cross arrow, or every outgoing
cross arrow, at either object. An extra internal arrow can
violate at most two conditions, and there are at most $s-6$
such arrows. An external violation requires a nonsingleton
hom-set and affects only one of the four conditions, so
the total number of violated conditions is at most
\[
 2(s-6)+P_{\rm ext}\leq s-4\leq3.
\]
Choose a condition that holds. The minimal arrows to and from
the other object have composite $e_i$ at the chosen object.
The unit hypothesis of Lemma~\ref{C:one-sided-corner-deletion}
holds because $e_i$ is a zero: for a unit $h$ and an incoming
$g$ fixed by $e_i$, one has $hg=he_i g=e_i g=g$; the outgoing
version is dual. That lemma deletes the object by a homotopy
equivalence. Its corner contributes at least two deleted
arrows, the other object contributes an arrow in each direction,
and nonextremality contributes an external entrance and an
external exit. At least six arrows disappear, leaving at most
fifteen, again impossible.

For $s=8$ the budget forces $P_{\rm ext}=0$.
The minimal group is trivial, and the possible total row and
column counts are $(2,2)$, $(3,2)$, or $(2,3)$.
There is therefore one minimal row per object or one minimal
column per object. All external hom-sets are singletons, so
all external entrances, respectively exits, are fixed by the
local minimal idempotents. Lemma~\ref{C:one-row-fixed-entrance-collapse}
or its dual replaces $S$ by a point through an actual
homotopy equivalence. Its eight internal arrows become one,
and no other hom-set increases. The result has at most
fourteen arrows. This exhausts the possibilities.
\end{proof}

\begin{theorem}\label{C:all-category-excess-eight-through21}
Every finite category with $k$ objects and
$n\leq\min\{21,2k+7\}$ arrows has classifying space not
homotopy equivalent to $\mathbb CP^2$.
\end{theorem}
\begin{proof}
Assume otherwise; the category is connected. Delete retract
objects successively. Each deletion is an actual homotopy
equivalence and removes at least three arrows (the deleted
identity and the two retract arrows). If $k'$ and $n'$ are
the remaining counts, this preserves
$n'\leq21$ and $n'\leq2k'+7$.

Suppose the remaining category has a proper strong component
on $t\geq3$ objects. Every identity of this component lies
outside its minimal ideal, so its internal arrow count is
at least $t^2+t$. There are $\ell=k'-t\geq1$ outside objects.
The connected support after contracting the component has
at least $\ell$ edges, including an interface which expands
to arrows at all $t$ objects of the component. With the
$\ell$ outside identities, the total is at least
\[
 n'\geq t^2+t+2\ell+(t-1)
       =2k'+t^2-1\geq2k'+8,
\]
a contradiction. A wholly strongly connected category is
excluded by the exact minimum twenty-three in
Corollary~\ref{C:strong-category-CP2-minimum23}.
A proper strong pair is excluded by the preceding proposition.
Otherwise every strong component is a single object and the
support is acyclic. Theorem~\ref{C:acyclic-arrow-excess-eight}
then excludes the remaining category. All cases are impossible.
\end{proof}

\begin{corollary}\label{C:CP2-object-bounds-through21}
A category of type $\mathbb CP^2$ with seventeen arrows has
at most four objects. With eighteen or nineteen arrows it
has at most five objects; with twenty or twenty-one arrows
it has at most six objects.
\end{corollary}
\begin{proof}
Apply the preceding theorem with, respectively,
$k\geq5$, $k\geq6$, and $k\geq7$.
\end{proof}
